% 自动生成：评测基准比较表（L4 层阅读卡，清洗版）
\begin{table*}[t]
\centering\scriptsize
\caption{评测基准比较（依据阅读卡证据整理，待人工复核；完整定位见证据库）}
\label{tab:benchmarks}
\begin{tabular}{llll}
\toprule
基准 & 年份 & 评测设置 & 主要发现（摘要） \\
\midrule
W0005 & 2024 & 使用 LOCOMO 基准，包含约 16K tokens、588.2 turns、27.2 sessions 的… & 作者称即使使用长上下文 LLM 或 RAG，模型在非常长期对话记忆上的表现仍远落后于人类，且在时间与因果动态理解上尤为困难… \\
W0006 & 2024 & 基准为 LongMemEval，包含 500 个手工问题，嵌入可扩展的用户-助手聊天历史… & 作者称此前长期记忆基准（如 LoCoMo）未评测回忆助手提供的信息或对更新后的用户信息进行推理 … \\
W0031 & 2026 & 基于 MERIT 基准，共评测 23,440 个得分 episode，总成本 \$42.57… & 总体：长期记忆将依赖任务成功率从无记忆的 0.00 提升至 0.55–1.00（pdf:p1 Abstract）… \\
W0171 & 2026 & Memora 基准包含多会话对话、基于记忆的评估问题和以每个 persona 显式记忆轨迹为依据的评测标准… & 作者称 Memora 在记忆整合和记忆变更要求上显著高于已有基准；作者称 LoCoMo 中 94\% 的评估问题只需不超过 2 个此前会… \\
W0379 & 2026 & 任务为跨会话多模态对话记忆评测；数据为 Mem-Gallery 数据集；模型为 13 个记忆机制… & 作者称基于 13 个记忆机制评测得出四项发现：(1) 显式保存视觉信息有益；(2) 记忆组织是必要的；(3) 推理与知识管理存在局限… \\
W0434 & 2026 & （证据未覆盖，待补读） & （证据未覆盖，待补读） \\
W0691 & 2025 & 提出MemBench基准，包含事实记忆与反思记忆两个层级，以及参与和观察两种交互场景… & 作者称MemBench的核心设计特征为同时覆盖记忆的两个层级与两种交互场景… \\
\bottomrule
\end{tabular}
\end{table*}
